% Generated draft: requires mathematical and source review.
\documentclass[11pt]{article}
\usepackage[margin=25mm]{geometry}
\usepackage{fontspec}
\setmainfont{Latin Modern Roman}
\usepackage{amsmath,amssymb,mathtools,unicode-math}
\setmathfont{Latin Modern Math}
\usepackage{enumitem,needspace,xurl,hyperref}
\setlength{\parindent}{0pt}
\setlength{\parskip}{5pt}
\setlength{\emergencystretch}{3em}
\newcommand{\N}{\mathbb{N}}
\newcommand{\Z}{\mathbb{Z}}
\newcommand{\Q}{\mathbb{Q}}
\newcommand{\R}{\mathbb{R}}
\newcommand{\C}{\mathbb{C}}
\newcommand{\F}{\mathbb{F}}
\newcommand{\A}{\mathbb{A}}
\newcommand{\PP}{\mathbb{P}}
\newcommand{\E}{\mathbb{E}}
\newcommand{\eps}{\varepsilon}
\newcommand{\dd}{\,\mathrm{d}}
\begin{document}
\section*{UnsolvedMath ID: 30003952}
\textbf{Source problem code:} catalogue entry
\subsection{Definitions and mathematical statement}
The supplied record does not specify the ambient surface, its boundary or marked points, or the conventions for homeomorphisms and isotopies. Consequently, the following is a conditional formulation requiring those choices before it becomes a fully specified problem.

Let \(S\) be the intended oriented surface, and let \(\mathcal H(S)\) denote its group of orientation-preserving homeomorphisms, subject to the intended boundary and marked-point conditions. Let \(\operatorname{Mod}(S)\) be the corresponding group of isotopy classes, with quotient map
\[
\pi\colon\mathcal H(S)\longrightarrow\operatorname{Mod}(S).
\]
Here isotopies must obey the same prescribed conditions. Let \(\alpha,\beta\subset S\) be essential simple closed curves meeting transversely in exactly one point. Essentiality must be interpreted according to the intended surface conventions, which the record does not supply. A right Dehn twist is the mapping class represented by one positive twist in an annular neighborhood of the curve, extended by the identity outside that neighborhood. Write
\[
a=T_\alpha,\qquad b=T_\beta,\qquad aba=bab.
\]
The homeomorphism realization question is whether there exist \(A,B\in\mathcal H(S)\) such that
\[
\pi(A)=a,\qquad \pi(B)=b,\qquad
A\circ B\circ A=B\circ A\circ B.
\]
The last equality is required as an equality of homeomorphisms.

If \(S\) is equipped with a compatible smooth structure, the regularity questions ask whether such representatives can be chosen to be \(C^1\) diffeomorphisms, or, for a prescribed positive integer \(r\), \(C^r\) diffeomorphisms, with the same boundary and marked-point conditions. The source does not specify the range or quantification of \(r\), nor whether its phrase \(C^r\) lifts requires \(C^r\) inverses. The diffeomorphism interpretation therefore requires confirmation. It also does not specify whether the question concerns a particular surface and curve pair or every admissible surface and pair.
\subsection{Short English statement}
Can two right Dehn twists about essential curves meeting once be represented by actual homeomorphisms satisfying the braid relation exactly? Can the representatives have \(C^1\) or \(C^r\) regularity? The ambient surface, boundary and marked-point conventions, quantifiers, and precise smoothness requirement remain unspecified. The supplied open-status assessment is dated curation and has not been independently verified.
\subsection{Sources}
\begin{itemize}
\item Local UnsolvedMath registry, permanent ID \(30003952\), source problem code \texttt{catalogue entry}, title: \emph{Smooth Lifts of the Dehn-Twist Braid Relation}; collection: Oberwolfach Reports Open Problems. No individual proposer is supplied. The record gives proposed year 2018.
\item Supplied source citation: \emph{New Trends in Teichm\"uller Theory and Mapping Class Groups} (2019), Oberwolfach Reports, pp.\ 2475--2534, DOI \texttt{10.4171/owr/2018/40}. The literature-source label identifies Oberwolfach Report 40, Problem 9. \url{https://doi.org/10.4171/owr/2018/40}.
\item Supplied literature source: Chen, non-realizability of braid groups. \url{https://arxiv.org/abs/1808.08248}.
\item Supplied literature source: Salter--Tshishiku, braid groups and diffeomorphisms. \url{https://doi.org/10.1112/blms/bdw016}.
\item Literature assessment supplied as checked on \texttt{2026-08-21}: the record labels the problem open and states that non-realization results for full braid or mapping class groups do not settle this two-generator realization question, which it describes as unresolved for homeomorphisms. This is supplied dated curation derived from the OWR record, not an independently verified literature assessment; its status and cited claims require independent checking.
\item The registry labels the statement \texttt{corrected\_verified} and reports that notation, typography, and omitted definitions were repaired while preserving mathematical content. These are supplied editorial claims, not mathematical verification of this draft.
\end{itemize}
\end{document}
